% HD-G23. GPT-6 Astra Pro, /root/g34; 21 November 2026, inherited effort.
\section{HD-G23: actual three-layer graph complexes and their trace kernel}
\label{sec:three-layer-graphs}

\paragraph{The actual finite construction.}
Use the six packages \(F_{J,\epsilon}=M_J\otimes\mathcal O_{B_\epsilon}\),
the integers \(c_\epsilon\), the ample line bundle \(L\), and
the negative sum \(N\) of HD-N16. Move the entire positive
\(B_{e_2}\) summand, in all six packages, to a third layer:
\[
 Q=\bigoplus_{|J|=2}F_{J,e_2}^{\,8(d+1)^3},\qquad
 P=L\oplus\bigoplus_{J,\ c_\epsilon>0,\ \epsilon\ne e_2}
                               F_{J,\epsilon}^{\,c_\epsilon}.
\]
First form an actual cone
\begin{equation}\label{eq:G23-F}
 N[1]\longrightarrow F\longrightarrow P
                              \xrightarrow{\delta}N[2]
\end{equation}
for any \(\delta\in\operatorname{Ext}^2(P,N)\). For every
actual morphism \(\beta:F\to Q[3]\), form the second actual cone
\begin{equation}\label{eq:G23-E}
 Q[2]\xrightarrow{i}E\xrightarrow{q}F
                                     \xrightarrow{\beta}Q[3].
\end{equation}
These are perfect complexes with three nonzero standard cohomology
layers \(Q,N,P\), in degrees \(-2,-1,0\). They have no negative
self-Ext, and their whole character is
\[
              \operatorname{ch}(E)=e^\theta+96w_d.
\]
There are choices with nonzero consecutive graph couplings and
nonzero coupling from \(L\), constructed explicitly below. For every
\(\delta\) and every actual \(\beta\), including all simultaneous
couplings which define such a morphism, the degree-two trace has a
kernel of dimension at least \(48\,104\,056\). Thus adding this
third layer does not give a semiregular object.

The assertion about transport is restricted to actual
Fourier--Mukai equivalences and isogeny pullbacks with their canonical
inherited deck actions, their subgroups, and actual equivalence
conjugates. A newly chosen coherent action after isogeny is outside
this dimension calculation.

\paragraph{Support separation and the negative Ext calculation.}
The positive point \(e_2=(0,1,0,0)\) is isolated in the graph
whose edges join positive slopes differing in one coordinate.
Indeed changing its second coordinate to zero gives the zero
coefficient, changing it to \(-1\) gives a negative coefficient,
and changing another coordinate to \(1\) or \(-1\) gives,
respectively, a negative double-positive coefficient or zero.
This follows directly from the coefficient formula in HD-G21:
\[
 c_{e_2}=8(d+1)^3,\quad c_{-e_2}=-8(d+1)^3,\quad
 c_{e_2+e_j}=-8d(d+1)^2,\quad c_{e_2-e_j}=0.
\]
Every other positive support therefore differs from \(B_{e_2}\)
in at least two elliptic coordinates. Distinct slopes in one
elliptic pair have local Ext only in degree one, so the four-factor
formula of HD-N16 gives
\[
 \operatorname{Hom}(P-L,Q)=\operatorname{Ext}^1(P-L,Q)=0.
\]
The line-to-graph cohomology has minimum degree two, whence also
\(\operatorname{Hom}(L,Q)=\operatorname{Ext}^1(L,Q)=0\).
All negative-positive support Hom groups vanish. In particular
\[
 \operatorname{Hom}(P,N)=\operatorname{Hom}(N,P)
     =\operatorname{Hom}(Q,N)=0,
 \qquad \operatorname{Ext}^{0,1}(P,Q)=0.
\]
The first cone \eqref{eq:G23-F} has no negative self-Ext by the
exact negative-degree formula of HD-G21.

For \(k<0\), the two Hom sequences of \eqref{eq:G23-F} show
\[
 \operatorname{Hom}(F,Q[2+k])=0,
 \qquad \operatorname{Hom}(Q[2],F[k])=0.
\]
For the first, the bounding groups are
\(\operatorname{Ext}^{k+2}(P,Q)\) and
\(\operatorname{Ext}^{k+1}(N,Q)\); the displayed vanishings
handle \(k=-1,-2\), and negative Ext handles the rest.
For the second, the bounding groups are
\(\operatorname{Ext}^{k-1}(Q,N)\) and
\(\operatorname{Ext}^{k-2}(Q,P)\), both negative.
Applying the two Hom sequences of \eqref{eq:G23-E}, with these
cross vanishings and the negative self-Ext vanishings of its two
pieces, proves \(\operatorname{Ext}^{<0}(E,E)=0\).
The standard cohomology sequence gives the three layers stated
above. Their alternating signs are \(+,-,+\), so moving \(Q\)
from degree zero to degree minus two preserves exactly the whole
character of HD-N16.

\paragraph{An explicit coherent three-layer coupling.}
Fix package \(J=\{3,4\}\), and choose one copy of each of
\[
 U=F_{J,e_1}\subset P,\qquad
 V=F_{J,e_1+e_2}\subset N,\qquad
 W=F_{J,e_2}\subset Q.
\]
The first pair of graphs differs only in elliptic coordinate two,
and the second pair only in coordinate one. Each transverse pair
has a nonzero degree-one Koszul Ext class. Let \(\xi_2\) and
\(\xi_1\) be those two classes. Let
\(\eta_3\in H^1(\mathcal O_{E_i})\) be a nonzero scalar class
on the shared third graph factor. The actual tensor-product Ext maps
\[
        u=\xi_2\otimes\eta_3\in\operatorname{Ext}^2(U,V),
 \qquad v=\xi_1\otimes\eta_3\in\operatorname{Ext}^2(V,W)
\]
are nonzero by Kunneth. Their Yoneda composite is zero, since its
third factor is \(\eta_3^2\in H^2(\mathcal O_{E_i})=0\).
This uses the actual common-factor tensor decomposition, not just
a count of dimensions.

There is also a nonzero
\(\ell\in\operatorname{Ext}^2(L,V)\): the line ratio on
\(B_{e_1+e_2}\) is negative in coordinates three and four and
trivial in coordinates one and two. Choose the tensor of nonzero
degree-one negative-line cohomology classes in coordinates three
and four, and the constant sections in the other two. Multiplication
by \(\eta_3\) kills this class, because an elliptic curve has no
degree-two coherent cohomology. Hence \(v\ell=0\) as well.

Define \(\delta\) by the two entries \(u\) and \(\ell\)
into the chosen \(V\), with its other entries zero, and let
\(\nu:N\to Q[2]\) have the single entry \(v\). We have proved
the literal Yoneda equation
\(\nu[2]\delta=0\in\operatorname{Ext}^4(P,Q)\).
Applying \(\operatorname{Hom}(-,Q[3])\) to
\eqref{eq:G23-F} gives a lift
\[
 \beta:F\longrightarrow Q[3],\qquad
 \beta|_{N[1]}=\nu[1].
\]
Thus both triangles are actual nonsplit cones, with nonzero
consecutive graph maps and nonzero line coupling.

The enhancement can be made explicit. Represent the three underlying
sheaf objects by complexes, choose cocycle representatives for
\(\delta\) and \(\nu\), and choose a degree-three cochain
\(h:P\to Q\) whose differential is \(-\nu\delta\), which
exists because the just-proved degree-four class is zero.
On the three shifted blocks \(P,N[1],Q[2]\), the upper triangular
differential has entries \(\delta,\nu,h\), with the usual shift
signs. Its only length-two off-diagonal square is
\(dh+\nu\delta=0\). There is no further block in which a longer
product could land. This produces the actual twisted complex
represented by the two cones; no higher compatibility is assumed.

\paragraph{An exact trace-kernel lift through the third layer.}
The two-layer triangle \eqref{eq:G23-F} has the actual trace-zero
map of HD-N16,
\[
 \operatorname{Ext}^3(P,N)\longrightarrow\operatorname{Ext}^2(F,F),
 \qquad a\longmapsto i_N[2]a q_P.
\]
Its kernel is exactly
\(\operatorname{Ext}^1(N,N)\delta+
\delta\operatorname{Ext}^1(P,P)\), by the two Hom sequences
and \(\operatorname{Hom}(N,P)=0\). Write \(T_F\) for its
image. Cyclic trace gives \(\sigma_F^2(T_F)=0\), and
\[
 \dim T_F\ge
  \operatorname{ext}^3(P,N)-\operatorname{ext}^1(P,P)
                                  -\operatorname{ext}^1(N,N).
\]
For the actual connecting map \(\beta\), restrict to
\[
 T_F^0=\ker\bigl(T_F\longrightarrow\operatorname{Ext}^5(F,Q),
                                   \ f\longmapsto\beta[2]f\bigr).
\]
From \eqref{eq:G23-F},
\(\operatorname{ext}^5(F,Q)
\le\operatorname{ext}^5(P,Q)+\operatorname{ext}^4(N,Q)\).
Therefore
\begin{equation}\label{eq:G23-bound}
 \dim T_F^0\ge R(d):=
 \operatorname{ext}^3(P,N)-\operatorname{ext}^1(P,P)
 -\operatorname{ext}^1(N,N)-\operatorname{ext}^5(P,Q)
 -\operatorname{ext}^4(N,Q).
\end{equation}

For each \(f\in T_F^0\), the exact Hom sequence of
\eqref{eq:G23-E} supplies \(\alpha:F\to E[2]\) with
\(q[2]\alpha=f\). Put \(z_f=\alpha q\).
Independent \(f\)'s give independent \(z_f\)'s for any choices
of lifts. Indeed \(\alpha q=0\) implies
\(\alpha=b\beta\) for some \(b:Q[3]\to E[2]\), by another
Hom sequence. But
\[
 q[2]b\in\operatorname{Hom}(Q[3],F[2])
          =\operatorname{Hom}(Q,F[-1])=0.
\]
The last vanishing follows from \eqref{eq:G23-F} shifted by minus
one: its bounding groups are \(\operatorname{Hom}(Q,N)=0\)
and \(\operatorname{Hom}(Q,P[-1])=0\). Thus \(f=0\), proving
independence. Finally cyclicity and Atiyah naturality give the exact
trace identity
\[
               \sigma_E^2(\alpha q)=\sigma_F^2(q[2]\alpha)
                                      =\sigma_F^2(f)=0.
\]
This proves a trace-kernel lower bound \(R(d)\) for every actual
\(\delta\) and \(\beta\), without assuming a generic rank or
discarding any connecting map.

\paragraph{Exact positive dimension polynomial.}
Every dimension in \eqref{eq:G23-bound} is a finite sum of the
four-factor Ext polynomials (N16.8), with precisely the package
sets \(P-L,N,Q\) specified above. The line-to-graph contributions
have degrees \(2,3,4\), so \(L\) contributes twice its degree-two
dimension to \(\operatorname{ext}^3(P,N)\), zero to
\(\operatorname{ext}^5(P,Q)\), and eight to
\(\operatorname{ext}^1(P,P)\). All other degree-one cross terms
with \(L\) vanish. The exact expansion, in \(t=d-1\), is
\begin{align*}
 R(1+t)={}&48\,104\,056+224\,295\,360t+452\,564\,448t^2
                      +517\,114\,896t^3\\
 &+367\,589\,952t^4+168\,088\,656t^5+49\,391\,904t^6\\
 &+9\,034\,560t^7+948\,768t^8+44\,928t^9.
\end{align*}
Every coefficient is positive. Hence
\(\dim\ker\sigma_E^2\ge48\,104\,056\) for every integer
\(d\ge1\). The exact program
\texttt{three\_layer\_single\_Q\_}\allowbreak\texttt{polynomial.py}
reconstructs the grid coefficients in \(\mathbb Z[t,s]/(s^2+t+1)\),
checks their uniform signs, multiplies the factor Ext polynomials,
and computes all five terms of \eqref{eq:G23-bound}. Its output
also retains the weaker exploratory bounds with their signed values;
only the positive quotient-lift polynomial just displayed is used
in this conclusion.

\paragraph{Point stabilizer and inherited equivariant trace.}
The standard cohomology layers are intrinsic. Three positive
single-coordinate graphs \(B_{e_1},B_{e_3},B_{e_4}\) occur in
degree zero, and \(B_{e_2}\) occurs in degree minus two. A point
functor preserving \(E\) must preserve each support slope in its
own degree. Their intersection is the origin by the same explicit
coordinate equations as HD-G21. Its degree-zero rank-one quotient
is \(L\), so the tensor coordinate is trivial as well.

Actual Fourier--Mukai equivalences preserve the Ext classes and
their categorical trace and identify the point functors by the
proved graph map of HD-G19. For a finite etale isogeny, self-Ext
pullback has the explicit adjunction/sheaf-trace left inverse and
preserves the zero trace. For the canonical deck action and each
of its subgroups, the actual descended object on the intermediate
quotient is again the pullback of \(E\); the same injective pullback
there proves the nonzero equivariant trace kernel. This is the
precise inherited-action argument of HD-N16. No conclusion about
a newly chosen post-isogeny linearization is added.

\paragraph{Scope of this completed three-layer test.}
Both adjacent couplings and the line coupling occur in actual
nonsplit complexes, and the third layer has positive dimension.
The proof covers every actual morphism \(\beta\) and every first
cone \(\delta\) with this fixed distribution of sheaf summands.
It does not assert that every other distribution of positive graphs
or every longer complex has the same bound.
